% ===== SECTION 4: OPERATIONAL EFFICIENCIES =====
\section{\texorpdfstring{\color{MidnightBlue}Operational Efficiencies}{Operational Efficiencies}}

\begin{sidebar}{RFI Reference}
Section C: Questions to Field --- Operational Efficiencies (Questions 1--3)
\end{sidebar}

% ----- Question 1 -----
\subsection{Question 1: Deadhead Reduction Strategies}
\textit{``In general, please describe your company's deadhead reduction strategies.''}

\vspace{0.3cm}

\placeholder{This is a critical question for RIDE given the extreme deadhead inefficiency in Zones C and D (2h12m to 3h38m of deadhead overhead per route). Sentry's response should demonstrate deep understanding of the problem and concrete technology-driven solutions.}

\vspace{0.3cm}

\mybox{
\textbf{\faIcon{exclamation-triangle} The Deadhead Problem in Context:} RIDE's own data reveals significant variation in deadhead time across zones. Zones C and D show average route times of 6h07m and 7h05m \textit{including} deadhead, but only 3h55m and 3h27m of productive service time---meaning 36--51\% of total route time is non-revenue deadhead mileage. With ridership up 34\% since 2015 (2,850 to 3,812 students) and close to 100 more buses per day on the road, deadhead reduction is now a \$41M system's largest efficiency lever.
}{MidnightBlue}

\vspace{0.3cm}

\placeholder{Describe Sentry's approach to deadhead reduction. Key strategies to cover:}

\begin{description}[style=nextline, leftmargin=0.5cm]
    \item[\faIcon{route} \textbf{Algorithm-Based Route Optimization:}] \placeholder{Describe how Sentry's routing engine minimizes deadhead by clustering pickups/drop-offs geographically, optimizing depot placement, and sequencing routes to minimize empty miles between runs.}

    \item[\faIcon{map-marked-alt} \textbf{Dynamic Depot Assignment:}] \placeholder{Describe how routes can be assigned to the nearest available depot or staging area rather than a fixed home base, reducing the initial deadhead leg.}

    \item[\faIcon{exchange-alt} \textbf{Multi-Run Chaining:}] \placeholder{Describe how AM and PM routes can be chained with midday routes (PCCT, McKinney-Vento) to keep vehicles productive and reduce return-to-depot deadhead between runs.}

    \item[\faIcon{border-all} \textbf{Cross-Zone Optimization:}] \placeholder{Describe how a centralized routing platform can identify opportunities to assign routes near zone boundaries to the closest depot regardless of zone, reducing deadhead that results from artificial geographic boundaries. Reference the Meeting Street consolidation example: centralized routing reduced buses serving that school from 28 to 13 (54\% reduction)---this same logic applied system-wide across 500+ destination schools is where transformative savings occur.}
\end{description}

\begin{sidebar}{Deadhead Reduction Potential}
\placeholder{Quantify the potential savings. Example: ``Based on RIDE's published route data, optimization of Zones C and D alone could reduce deadhead time by an estimated 25--35\%, translating to approximately X fewer bus-hours per day and \$X in annual fuel and labor savings.''}
\end{sidebar}

% ----- Question 2 -----
\subsection{Question 2: Preventive Maintenance Efficiency Metrics}
\textit{``Please describe your company's preventive maintenance efficiency metrics and how those relate to how your organization best-manages its preventative maintenance program.''}

\vspace{0.3cm}

\placeholder{Describe Sentry's approach to preventive maintenance tracking and fleet management. Sentry provides the \textit{technology platform} that enables operators to manage PM programs---not the wrench-turning itself.}

\begin{description}[style=nextline, leftmargin=0.5cm]
    \item[\faIcon{tools} \textbf{PM Scheduling Engine:}] \placeholder{Describe automated maintenance scheduling based on mileage, engine hours, calendar intervals, and manufacturer specifications. System generates work orders and tracks completion.}

    \item[\faIcon{chart-bar} \textbf{Fleet Health Dashboard:}] \placeholder{Describe real-time visibility into fleet status---vehicles due for PM, overdue units, out-of-service rates, mean time between failures (MTBF), and parts inventory levels.}

    \item[\faIcon{clipboard-check} \textbf{Pre- \& Post-Trip Inspection Digitization:}] \placeholder{Describe how drivers complete inspections on tablets/phones, defects are flagged immediately, and critical defects trigger automatic vehicle holds until resolved.}

    \item[\faIcon{tachometer-alt} \textbf{Key Metrics Tracked:}]
    \begin{itemize}
        \item \textbf{PM Compliance Rate:} \placeholder{Target \% and how it's measured}
        \item \textbf{Mean Distance Between Failures (MDBF):} \placeholder{Target and benchmark}
        \item \textbf{Spare Ratio Utilization:} \placeholder{Target \% of spares deployed on any given day}
        \item \textbf{Road Call Rate:} \placeholder{Target per 100,000 miles}
    \end{itemize}
\end{description}

% ----- Question 3 -----
\subsection{Question 3: Driver Scheduling Efficiency Tools and Metrics}
\textit{``Please describe your company's driver scheduling efficiency tools and metrics.''}

\vspace{0.3cm}

\placeholder{Describe Sentry's driver scheduling capabilities. Emphasize how technology reduces overtime, improves driver retention, and ensures compliance with hours-of-service rules.}

\begin{description}[style=nextline, leftmargin=0.5cm]
    \item[\faIcon{calendar-alt} \textbf{Automated Bid/Assignment System:}] \placeholder{Describe how drivers are matched to routes based on qualifications (CDL endorsements, special education training, language skills), seniority (union compliance), and geographic proximity to reduce personal deadhead.}

    \item[\faIcon{user-clock} \textbf{Hours Optimization:}] \placeholder{Describe how the system balances driver hours across the workforce to minimize overtime while ensuring adequate rest periods and compliance with DOT regulations.}

    \item[\faIcon{bell} \textbf{Real-Time Absence Management:}] \placeholder{Describe automated call-out detection, standby driver notification, and route reassignment to minimize service disruptions from driver absences.}

    \item[\faIcon{chart-pie} \textbf{Key Metrics:}]
    \begin{itemize}
        \item \textbf{Driver Utilization Rate:} \placeholder{Productive hours as \% of paid hours}
        \item \textbf{Overtime as \% of Total Hours:} \placeholder{Target and benchmark}
        \item \textbf{Open Route Rate:} \placeholder{Routes without assigned driver on any given day}
        \item \textbf{Driver Turnover Rate:} \placeholder{Annual \% and improvement strategies}
    \end{itemize}
\end{description}
